\begin{thebibliography}{99}
\footnotesize\itemsep=0pt

\bibitem{AbMa}
Abraham R., Marsden J.E., Foundations of mechanics, 2nd~ed., Benjamin/Cummings
 Publishing Co., Inc., Advanced Book Program, Reading, Mass., 1978.

\bibitem{Arn}
Arnold V., Les m\'ethodes math\'ematiques de la m\'ecanique classique, Mir,
 Moscow, 1976.

\bibitem{BouBr}
Boualem H., Brouzet R., Topology of the space of locally separable functions,
 \href{https://doi.org/10.1016/j.topol.2021.107728}{\textit{Topology Appl.}} \textbf{299} (2021), 107728, 16~pages.

\bibitem{Br}
Brouzet R., About the existence of recursion operators for completely
 integrable {H}amiltonian systems near a {L}iouville torus, \href{https://doi.org/10.1063/1.530212}{\textit{J.~Math.
 Phys.}} \textbf{34} (1993), 1309--1313.

\bibitem{BrMoTu}
Brouzet R., Molino P., Turiel F.J., G\'eom\'etrie des syst\`emes
 bihamiltoniens, \href{https://doi.org/10.1016/0019-3577(93)90002-G}{\textit{Indag. Math.~(N.S.)}} \textbf{4} (1993), 269--296.

\bibitem{AMM1}
Craik A.D.D., Prehistory of {F}a\`a di {B}runo's formula, \href{https://doi.org/10.2307/30037410}{\textit{Amer. Math.
 Monthly}} \textbf{112} (2005), 119--130.

\bibitem{Ehr}
Ehresmann C., Les connexions infinit\'esimales dans un espace fibr\'e
 diff\'erentiable, in Colloque de topologie (espaces fibr\'es), {B}ruxelles,
 1950, Georges Thone, Li\`ege; Masson \& Cie, Paris, 1951, 29--55.

\bibitem{FP}
Falqui G., Pedroni M., Separation of variables for bi-{H}amiltonian systems,
 \href{https://doi.org/10.1023/A:1024080315471}{\textit{Math. Phys. Anal. Geom.}} \textbf{6} (2003), 139--179,
 \href{https://arxiv.org/abs/nlin.SI/0204029}{arXiv:nlin.SI/0204029}.

\bibitem{Fe}
Fernandes R.L., Completely integrable bi-{H}amiltonian systems,
 \href{https://doi.org/10.1007/BF02219188}{\textit{J.~Dynam. Differential Equations}} \textbf{6} (1994), 53--69.

\bibitem{GD}
Gel'fand I.M., Dorfman I.Ja., Hamiltonian operators and algebraic structures
 related to them, \href{https://doi.org/10.1007/BF01078363}{\textit{Funct. Anal. Appl.}} \textbf{13} (1979), 248--262.

\bibitem{Har}
Hardy M., Combinatorics of partial derivatives, \href{https://doi.org/10.37236/1027}{\textit{Electron.~J. Combin.}}
 \textbf{13} (2006), 1, 13~pages, \href{https://arxiv.org/abs/math.CO/0601149}{arXiv:math.CO/0601149}.

\bibitem{Hi}
Hirsch M.W., Differential topology, \textit{Graduate Texts in Mathematics},
 Vol.~33, \href{https://doi.org/10.1007/978-1-4684-9449-5}{Springer-Verlag}, New York~-- Heidelberg, 1976.

\bibitem{IMM}
Ibort A., Magri F., Marmo G., Bihamiltonian structures and {S}t\"ackel
 separability, \href{https://doi.org/10.1016/S0393-0440(99)00051-0}{\textit{J.~Geom. Phys.}} \textbf{33} (2000), 210--228.

\bibitem{AMM2}
Johnson W.P., The curious history of {F}a\`a di {B}runo's formula,
 \href{https://doi.org/10.2307/2695352}{\textit{Amer. Math. Monthly}} \textbf{109} (2002), 217--234.

\bibitem{Mar}
Landi G., Marmo G., Vilasi G., Recursion operators: meaning and existence for
 completely integrable systems, \href{https://doi.org/10.1063/1.530613}{\textit{J.~Math. Phys.}} \textbf{35} (1994),
 808--815.

\bibitem{LiMa}
Libermann P., Marle C.-M., G\'eom\'etrie symplectique, bases th\'eoriques de la
 m\'ecanique. {T}ome {I}, \textit{Publications Math\'ematiques de
 l'Universit\'e Paris~VII}, Vol.~21, Universit\'e de Paris~VII, U.E.R. de
 Math\'ematiques, Paris, 1986.

\bibitem{Ma}
Ma T.-W., Higher chain formula proved by combinatorics, \href{https://doi.org/10.37236/259}{\textit{Electron.~J.
 Combin.}} \textbf{16} (2009), 21, 7~pages.

\bibitem{Mag}
Magri F., A simple model of the integrable {H}amiltonian equation,
 \href{https://doi.org/10.1063/1.523777}{\textit{J.~Math. Phys.}} \textbf{19} (1978), 1156--1162.

\bibitem{MaMo}
Magri F., Morosi C., A geometrical characterization of integrable {H}amiltonian
 systems through the theory of {P}oisson--{N}ijenhuis manifolds, {D}epartment
 of Mathematics, University of Milan, 1984, available at
 \url{https://boa.unimib.it/retrieve/handle/10281/17656/21540/quaderno3-2008.pdf}.

\bibitem{Nij}
Nijenhuis A., Jacobi-type identities for bilinear differential concomitants of
 certain tensor fields, \textit{Indag. Math.} \textbf{58} (1955), 390--403.

\bibitem{PS}
Praught J., Smirnov R.G., Andrew {L}enard: a mystery unraveled, \href{https://doi.org/10.3842/SIGMA.2005.005}{\textit{SIGMA}}
 \textbf{1} (2005), 005, 7~pages, \href{https://arxiv.org/abs/nlin.SI/0510055}{arXiv:nlin.SI/0510055}.

\bibitem{TT}
Tempesta P., Tondo G., Haantjes algebras of classical integrable systems,
 \href{https://doi.org/10.1007/s10231-021-01107-4}{\textit{Ann. Mat. Pura Appl.}}, {t}o appear, \href{https://arxiv.org/abs/1405.5118}{arXiv:1405.5118}.

\end{thebibliography}
